% !TeX program = xelatex
\documentclass[9pt]{beamer}
\usepackage[szu,en]{collegeBeamer}
\usepackage{fontspec}
% falls back to a widely-available monospace font on machines without Cascadia Code installed

\IfFontExistsTF{Cascadia Code}
    {\newfontfamily\cascadia{Cascadia Code}}
    {\newfontfamily\cascadia{DejaVu Sans Mono}}
\usetikzlibrary{positioning,arrows.meta}

% biblatex: reuses relevant entries from tesis/Ref/SampleReferences.bib (copied into
% referencias.bib inside presentation/, without touching the original .bib in tesis/)
\usepackage[backend=biber,style=numeric,sorting=none]{biblatex}
\addbibresource{referencias.bib}

\renewcommand{\qaStr}{\textsl{Thank you for your attention! \\ Questions and comments are welcome.}}

\definecolor{hrefcol}{RGB}{55, 50, 204} % Example: blue color

% meta-data
\title{
{\large\color{blue!60!black}\textbf{Computer Science}}\\[0.8em]
Extraction and longitudinal \\
monitoring of clinical information\\
in hepatology using GraphRAG
}
% the theme's title-page template always appends "\\" after \insertsubtitle;
% an empty subtitle makes that "\\" have nothing to end, causing a LaTeX error,
% so a near-invisible placeholder (thin space) is used instead of leaving it undefined
\subtitle{\,}

\author{
Matheus Ryan da Silva Nascimento\\
{\footnotesize\cascadia\href{mailto:mrysn@ic.ufal.br}{mrysn@ic.ufal.br}}}

\date{
Maceió, Brazil\\
August 2026
}

% document body
\begin{document}

    \maketitle

    \section{Introduction}

    \begin{frame}{Context and Motivation}
        \begin{itemize}
            \item Liver diseases: a major global public-health problem \\
            (cirrhosis, viral hepatitis, hepatocellular carcinoma) \cite{sousa2023}
            \item In Brazil, chronic hepatitis mortality has been rising in the North/Northeast (2001--2020) \cite{sousa2023}
            \item Diagnosis and follow-up rely on electronic health records rich in free text \cite{schneider2023}
            \item Lack of computational tools hinders early detection and treatment planning \cite{hou2014}
        \end{itemize}
    \end{frame}

    \begin{frame}{NLP, RAG, and GraphRAG}
        \small
        \begin{itemize}
            \itemsep1pt
            \item NLP has advanced in gastroenterology/hepatology, including with LLMs, but still faces abbreviations, clinical jargon, and temporal evolution \cite{stammers2025,shahab2024}
            \item Tools like ChatGPT help synthesize medical information, with limitations in accuracy and bias \cite{dave2023}
            \item \textbf{RAG}: retrieval integrated with LLM-based generation \cite{gao2024}
            \item \textbf{GraphRAG}: LLMs + knowledge graphs -- models hierarchies and complex relations between clinical entities \cite{edge2024}
            \item Effectiveness depends on the task, graph quality, and community-detection strategy \cite{xiang2025}
        \end{itemize}
        \vskip0.3em
        \centering\footnotesize\ctextbf{Gap: longitudinal analysis of clinical progression in liver-disease patients remains largely unexplored}
    \end{frame}

    \begin{frame}{General Objective}
        Develop and evaluate a \textbf{GraphRAG}-based approach for extracting and monitoring clinical information in electronic health records of liver-disease patients.

        \vskip1em
        \begin{block}{Research Questions}
            \begin{enumerate}
                \item How can context summarization improve the analysis of relevant clinical information?
                \item To what extent does temporal modeling with GraphRAG reveal meaningful evolutionary patterns in liver-disease progression?
                \item What are the limitations of integrating GraphRAG with unstructured clinical data (accuracy and relevance)?
            \end{enumerate}
        \end{block}
    \end{frame}

    \begin{frame}{Specific Objectives}
        \begin{itemize}
            \item Pre-process and anonymize clinical progress notes (named entity recognition)
            \item Structure the text base into semantic blocks (\textit{text chunks}) for analysis with LLMs
            \item Extract entities and relations to build the knowledge graph (KG)
            \item Compute statistical and descriptive metrics of the dataset (\texttt{TF-IDF}, \texttt{t-SNE})
            \item Develop prompts for summarization and generation of contextualized global answers
            \item Evaluate precision, relevance, and summarization quality, comparing segmentation and clustering strategies
        \end{itemize}
    \end{frame}

    \section{Relevance of the Work}

    \begin{frame}{Why This Work Matters}
        \small
        \begin{itemize}
            \itemsep1pt
            \item Liver diseases carry a heavy burden of hospitalizations and mortality; the literature points to the need for precision medicine in hepatology \cite{SABRIE2023350}
            \item Free-text records remain underused: NLP already extracts useful information, but with low explainability and little temporal integration \cite{omar2025}
            \item NLP + knowledge graphs + retrieval-augmented generation (GraphRAG) is still rare in hepatology \cite{stammers2025}
        \end{itemize}
        \vskip0.2em
        \begin{block}{Three gaps addressed}
            \begin{enumerate}
                \itemsep0pt
                \item Intelligent extraction of unstructured textual data in liver disease
                \item Longitudinal monitoring in a secure, local environment
                \item Use of knowledge graphs to explore this information
            \end{enumerate}
        \end{block}
    \end{frame}

    \section{Methodology}

    \begin{frame}{Pipeline Overview}
        \begin{center}
        \resizebox{\textwidth}{!}{%
        \begin{tikzpicture}[node distance=0.5cm, every node/.style={font=\scriptsize}]
            \tikzset{box/.style={draw=maincolor, thick, rounded corners, fill=morelightgray, text width=2.3cm, align=center, minimum height=1.5cm}}
            \node[box] (a) {Clinical progress notes};
            \node[box, right=of a] (b) {Pre-processing \& anonymization};
            \node[box, right=of b] (c) {Semantic chunking};
            \node[box, right=of c] (d) {Entity/relation extraction (KG)};
            \node[box, right=of d] (e) {Queries \& answer generation};
            \foreach \x/\y in {a/b, b/c, c/d, d/e}{\draw[-{Latex}, thick] (\x) -- (\y);}
        \end{tikzpicture}%
        }
        \end{center}
        \vskip0.5em
        Entity extraction and relation linking are handled as distinct steps, ensuring that complex, inter-related information is correctly represented in the KG.
    \end{frame}

    \begin{frame}{Dataset and Privacy}
        \begin{itemize}
            \item Clinical progress notes of liver-disease patients provided by UFAL's University Hospital (HU-UFAL)
            \item \textbf{100\% local} processing, via \texttt{graphrag-local-ollama}, using the \texttt{mistral} model
            \item Sensitive data never leaves the controlled environment -- compliant with Brazil's \textbf{LGPD}
            \item Anonymization via NER: \texttt{spaCy} framework + local \texttt{gpt-oss:8b} model (patient names, physicians, and other sensitive data)
            \item Final dataset split into training, validation, and test sets; documents segmented into \textit{chunks} with minimal overlap
        \end{itemize}
    \end{frame}

    \begin{frame}{Entity and Relation Extraction}
        \begin{itemize}
            \item Entities extracted by the LLM from text chunks: patients, exams, diagnoses, medications, procedures, clinical events
            \item Concise descriptions per entity (dates, exam values, interactions)
            \item Relations such as ``patient underwent exam'', ``diagnosis associated with exam'', ``medication prescribed for diagnosis''
            \item Repeated entities are aggregated; frequency weighs the relations
            \item \texttt{Distropic} embeddings for semantic similarity and consolidation across multiple chunks
        \end{itemize}
    \end{frame}

    \begin{frame}{Query and Information Retrieval}
        \begin{itemize}
            \item Community summaries organized into blocks, processed by the LLM with task-specific prompts
            \item The LLM assigns a relevance score to each intermediate answer, filtering out irrelevant/redundant results
            \item Selected answers are aggregated into a global context, respecting the token limit
            \item Query--text similarity computed via \texttt{Distropic} embeddings, prioritizing the most pertinent information
        \end{itemize}
    \end{frame}

    \section{Project Status}

    \begin{frame}{Planned Timeline}
        \begin{center}
        \small
        \renewcommand{\arraystretch}{1.15}
        \begin{tabular}{|p{5.6cm}|c|c|c|}
            \hline
            \textbf{Phase} & \textbf{Months 1--2} & \textbf{Months 3--4} & \textbf{Months 5--6} \\
            \hline
            Dataset collection \& anonymization & X & & \\
            \hline
            Semantic segmentation \& entity extraction & X & X & \\
            \hline
            Graph construction \& querying & & X & X \\
            \hline
            Evaluation, refinement, and final write-up & & & X \\
            \hline
        \end{tabular}
        \end{center}
        \vskip0.5em
        {\scriptsize Condensed from the project's 6-month plan; carried out incrementally, with continuous review.}
    \end{frame}

    \begin{frame}{Current Status}
        % TODO: update this list with the real status of each stage before the defense
        \begin{itemize}
            \item \ctextbf{Completed:} pre-processing, anonymization (NER), and semantic segmentation of the clinical-notes dataset
            \item \ctextbf{In progress:} entity/relation extraction and knowledge-graph construction; query and information retrieval
            \item \ctextbf{Next steps:} systematic evaluation (precision, relevance, summarization quality) comparing segmentation and clustering strategies
        \end{itemize}
    \end{frame}

    \section{Concluding Remarks}

    \begin{frame}{Work in Progress}
        % TODO: replace with real conclusions and results once the experiments are complete
        \begin{itemize}
            \item GraphRAG approach applied to a still largely unexplored domain: hepatology
            \item Local, secure pipeline, LGPD-compliant, from pre-processing to answer generation
            \item Comparative evaluation between segmentation and entity-clustering strategies under way
            \item Expected contribution: a replicable method for longitudinal monitoring of liver-disease patients from unstructured textual data
        \end{itemize}
    \end{frame}

    \bibliographpage

    \QApage

\end{document}
